% ============================================================
% CAPÍTULO 10: RUTAS METABÓLICAS CATABÓLICAS
% ============================================================
% CONTENIDO BASADO EN: Unidad de Aprendizaje III, Tema "Rutas metabólicas catabólicas"
% PROPÓSITO: Explicar las rutas catabólicas de las biomoléculas en los diferentes organismos vivos.
% ============================================================

\section{Introducción al catabolismo}

% --- CONSEJO: Define el catabolismo como la fase degradativa del metabolismo, que libera energía a partir de la ruptura de moléculas complejas en moléculas más simples. ---
\resumebox{definicion}{Catabolismo}{Es el conjunto de rutas metabólicas que degradan moléculas orgánicas complejas en moléculas más simples, liberando energía que se conserva en forma de ATP y poder reductor (NADH, FADH2).}

\subsection{Características generales del catabolismo}
\begin{itemize}
    \item \textbf{Fase de degradación:} Las biomoléculas (carbohidratos, lípidos, proteínas) se rompen en unidades más pequeñas.
    \item \textbf{Generación de energía:} Se produce ATP, NADH y FADH2.
    \item \textbf{Oxidación:} Implica reacciones de oxidación que liberan electrones.
    \item \textbf{Convergencia:} Las diferentes rutas catabólicas convergen en vías comunes (ej. ciclo del ácido cítrico).
\end{itemize}

\section{Glucólisis y metabolismo de hexosas y pentosas}

\subsection{Glucólisis}

% --- CONSEJO: Es la ruta central del catabolismo de carbohidratos. Ocurre en el citosol y degrada la glucosa a piruvato. ---
\resumebox{definicion}{Glucólisis}{Es la ruta metabólica que degrada una molécula de glucosa (6 carbonos) en dos moléculas de piruvato (3 carbonos), generando 2 ATP y 2 NADH en el proceso. No requiere oxígeno (es anaeróbica).}

% --- CONSEJO: Explica brevemente las dos fases de la glucólisis: la fase de inversión de energía (consumo de 2 ATP) y la fase de generación de energía (producción de 4 ATP y 2 NADH). ---

\subsection{Vía de las pentosas fosfato}

% --- CONSEJO: Ruta alternativa para la oxidación de la glucosa. Es importante para generar NADPH (poder reductor para la biosíntesis) y ribosa-5-fosfato (precursor de nucleótidos). ---
\resumebox{definicion}{Vía de las pentosas fosfato}{Es una ruta metabólica que oxida la glucosa-6-fosfato. Su principal función es generar NADPH y ribosa-5-fosfato. Ocurre en el citosol.}

\section{Ciclo del ácido cítrico (Ciclo de Krebs)}

% --- CONSEJO: Es la ruta final común para la oxidación de carbohidratos, lípidos y proteínas. Ocurre en la matriz mitocondrial (en eucariotas). ---
\resumebox{definicion}{Ciclo del ácido cítrico}{Es una ruta metabólica cíclica que oxida completamente el acetil-CoA a CO2, generando GTP, NADH y FADH2. Es una ruta anfibólica (participa en catabolismo y anabolismo).}

% --- CONSEJO: Explica que el acetil-CoA se condensa con oxaloacetato para formar citrato, y a través de una serie de reacciones se regenera el oxaloacetato, liberando CO2 y generando poder reductor. ---

\section{Cadena transportadora de electrones y fosforilación oxidativa}

% --- CONSEJO: Es la ruta final del catabolismo aeróbico. Ocurre en la membrana mitocondrial interna (en eucariotas). ---
\resumebox{definicion}{Cadena transportadora de electrones}{Es un conjunto de complejos proteicos en la membrana mitocondrial interna que transfieren electrones desde el NADH y FADH2 al O2, generando un gradiente de protones que se utiliza para sintetizar ATP.}

% --- CONSEJO: Explica el flujo de electrones (NADH -> complejo I -> Q -> complejo III -> cit c -> complejo IV -> O2), el bombeo de protones y la síntesis de ATP por la ATP sintasa. ---

\section{Fermentación}

% --- CONSEJO: Es una ruta anaeróbica que regenera el NAD+ a partir del NADH, permitiendo que la glucólisis continúe en ausencia de oxígeno. ---
\subsection{Fermentación alcohólica}
% --- CONSEJO: Ocurre en levaduras y algunas bacterias. El piruvato se descarboxila a acetaldehído, que luego se reduce a etanol, regenerando NAD+. ---

\subsection{Fermentación láctica}
% --- CONSEJO: Ocurre en algunas bacterias y en el músculo esquelético. El piruvato se reduce directamente a lactato, regenerando NAD+. ---

\section{Beta-oxidación de ácidos grasos}

% --- CONSEJO: Es la ruta principal para la degradación de ácidos grasos. Ocurre en la matriz mitocondrial. ---
\resumebox{definicion}{Beta-oxidación}{Es la ruta catabólica que degrada los ácidos grasos en unidades de acetil-CoA, eliminando dos carbonos a la vez. Genera NADH y FADH2.}

% --- CONSEJO: Explica el ciclo de cuatro reacciones que acorta la cadena del ácido graso en dos carbonos, liberando acetil-CoA. ---

\section{Oxidación de aminoácidos y ciclo de la urea}

% --- CONSEJO: El catabolismo de aminoácidos implica la eliminación del grupo amino (desaminación) y la oxidación del esqueleto carbonado. ---
\subsection{Desaminación y transaminación}
% --- CONSEJO: Explica cómo los grupos amino se transfieren al alfa-cetoglutarato para formar glutamato, que luego se desamina liberando amoniaco (NH3). ---

\subsection{Ciclo de la urea}
% --- CONSEJO: El amoniaco es tóxico para las células. El ciclo de la urea (en el hígado) convierte el amoniaco en urea, que es menos tóxica y se excreta. ---

\section{Metilotróficas}

% --- CONSEJO: Este es un tema más avanzado. Las vías metilotróficas permiten a algunos microorganismos utilizar compuestos de un solo carbono (como metanol, metano, CO2) como fuente de carbono y energía. ---
\resumebox{definicion}{Vías metilotróficas}{Son rutas metabólicas que permiten a ciertos microorganismos (metilotrofos) crecer utilizando compuestos de un solo carbono (como metanol o metano) como única fuente de carbono y energía.}

\section{Integración de las rutas catabólicas}

% --- CONSEJO: El objetivo final es mostrar cómo se integran todas estas rutas. La glucosa, los ácidos grasos y los aminoácidos se degradan a acetil-CoA, que entra al ciclo del ácido cítrico, generando poder reductor que se oxida en la cadena transportadora de electrones para producir ATP. ---